\begin{textsample}{POS Dim 1 – mg – Score 49.00 – mg/mg\_em\_71.txt}  \label{ex:f1_pos_234}
3.4 A área de Ciências Humanas e Sociais Aplicadas A Área de Conhecimento de Ciências Humanas e Sociais Aplicadas – constituída pelos componentes curriculares Filosofia, Geografia, História e Sociologia – integra o presente Currículo Referência tendo como \textbf{foco} principal propiciar aos estudantes a \textbf{compreensão} de si, enquanto sujeitos autônomos e protagonistas de sua formação escolar.
Para \textbf{atingir} tal propósito, \textbf{fomenta} uma dinâmica educacional voltada para a apreensão e o desenvolvimento de \textbf{competências} e habilidades \textbf{indispensáveis} à formação intelectual, ética e cidadã do sujeito.
Esse campo do conhecimento se fundamenta a partir de uma relação interdisciplinar de seus componentes curriculares, reconhecendo os direitos de aprendizagem e desenvolvimento dos estudantes ( BRASIL, 2018, p. 7 ), assim como promovendo a ampliação de suas visões de mundo.
Isso porque, romper com uma tradição escolar, que se vale da fragmentação, da compartimentalização e da hierarquização dos saberes, se faz necessário em nosso tempo.
Nesta perspectiva, o Currículo Referência, ao estimular a maior interação entre os componentes curriculares que compõem cada Área de Conhecimento, e, particularmente, a Área de Ciências Humanas e Sociais Aplicadas, intenta estimular nos professores que ministram as \textbf{aulas} de Filosofia, Geografia, História e Sociologia a busca, em conjunto, de estratégias que \textbf{sinalizem} > “ a superação da fragmentação radicalmente disciplinar do conhecimento, o estímulo à sua aplicação na vida real, a importância do contexto para dar sentido ao que se aprende e o protagonismo do estudante em sua aprendizagem e na construção de seu projeto de vida ”. ( BRASIL 2018, p. 15 ) A \textbf{interdisciplinaridade}, portanto, torna -se eixo fundamental nessa nova perspectiva educacional, pois pressupõe o diálogo entre os princípios \textbf{metodológicos} e epistemológicos dos componentes curriculares que compõem a Área de Ciências Humanas e Sociais Aplicadas.
Com relação a essa nova configuração para a \textbf{abordagem} dos objetos de conhecimento ao longo da etapa do Ensino Médio, a BNCC esclarece que o currículo deve > decidir sobre formas de organização interdisciplinar dos componentes curriculares e fortalecer a \textbf{competência} pedagógica das equipes escolares para adotar estratégias mais dinâmicas, interativas e colaborativas em relação à gestão do ensino e da aprendizagem. ( BRASIL, 2018, p. 16 ) Vale ressaltar que a junção da Filosofia, da Geografia, da História e da Sociologia em uma única Área do Conhecimento não representa a supressão de nenhum desses quatro componentes curriculares e de suas especificidades.
Na verdade, \textbf{sinaliza} para a potência de maior interação entre esses saberes na promoção de um efetivo \textbf{letramento} dos estudantes na etapa do Ensino Médio para o exercício da cidadania, o ingresso \textbf{qualificado} no mundo do trabalho, a proteção do meio ambiente, o incentivo a um modo de vida sustentável, a defesa dos direitos humanos, a fruição consciente das novas \textbf{tecnologias} da informação e da comunicação, o desenvolvimento de habilidades \textbf{socioemocionais} e a construção de um projeto de vida voltado para a ampliação do campo de possibilidade desses sujeitos.
Sendo a BNCC o documento normativo que define os \textbf{parâmetros} para a elaboração dos currículos escolares no que \textbf{diz} respeito à educação básica, é de \textbf{suma} importância nos atentarmos para a estruturação e a arquitetura \textbf{conceitual} de seu texto.
Essa nova perspectiva de educação que \textbf{inspira} e fundamenta o Currículo Referência se pauta por valores e conceitos humanistas - encorajando o protagonismo juvenil e a formação integral do sujeito.
Vale ressaltar que uma “ educação humanista consiste, antes de tudo, em \textbf{fomentar} e ensinar o uso da \textbf{razão}, essa \textbf{capacidade} que observa, abstrai, deduz, \textbf{argumenta} e conclui logicamente. ” ( SAVATER, 2005, p. 131 ).
Assim sendo, ao nos debruçarmos sobre os objetivos que orientam o desenvolvimento das Áreas do Conhecimento que compõem a BNCC, podemos apreciar, compreender melhor o \textbf{horizonte} humanista que fortalece essa nova proposta de formação educacional e comprometermo -nos com a sua implementação.
A Área de Linguagens e suas Tecnologias, por exemplo, enfatiza que experiências propiciadas por esse saber devem ser “ situadas em campos de atuação social diversos, vinculados com o enriquecimento cultural próprio, as práticas cidadãs, o trabalho e a continuação dos estudos ” ( BNCC, p. 485 ).
Por sua vez, a Área de Matemática e suas Tecnologias estabelece que, para uma real efetivação de seus propósitos, “ os estudantes devem desenvolver habilidades relativas aos processos de investigação, de construção de modelos e de resolução de \textbf{problemas} ” ( BNCC, p. 529 ).
Já com relação à Área de Ciências da Natureza e suas Tecnologias é ressaltado que “ a \textbf{contextualização} social, histórica e cultural da \textbf{ciência} e da \textbf{tecnologia} é fundamental para que elas sejam compreendidas como empreendimentos humanos e sociais ” ( BNCC, p. 549 ).
Por fim, a Área de Ciências Humanas e Sociais Aplicadas também apresenta características fundamentais para a nova configuração curricular.
Como preconiza a própria BNCC, os estudos dos componentes Filosofia, Geografia, História e Sociologia proporcionam um “ domínio maior sobre diferentes linguagens, o que favorece os processos de simbolização e de abstração ”, e “ propõe que os estudantes desenvolvam a \textbf{capacidade} de estabelecer diálogos ” ( Ibidem ), “ elaborem hipóteses e \textbf{argumentos} com base na seleção e na \textbf{sistematização} de dados, obtidos em fontes confiáveis e sólidas ” ( BRASIL, p. 561, 562 ), além de levá -los à prática da “ dúvida sistemática – entendida como questionamento e autoquestionamento, \textbf{conduta} contrária à crença em verdades absolutas ” ( Ibidem ).
Assim sendo, percebemos que, para além de promover e desenvolver as especificidades \textbf{inerentes} aos componentes curriculares Filosofia, Geografia, História e Sociologia – elementos esses que propiciam aos estudantes a \textbf{compreensão} do mundo enquanto processo, sob \textbf{influências} de determinismos naturais e sociais, em contínua construção e essencialmente aberto à intervenção humana ( o que por si só já justifica o inestimável valor desses quatro campos do saber presentes no Currículo Referência Estadual ) – a Área de Ciências Humanas e Sociais Aplicadas contribui para a consolidação dos aprendizados explorados na etapa do Ensino Fundamental e apresenta, em seu âmago, elementos fundamentais para um aprendizado mais significativo das outras três Áreas do Conhecimento que integram o Currículo Referência.
Em \textbf{suma}, “ aprender a discutir, a refutar e a justificar o que se pensa é parte \textbf{indispensável} de qualquer educação que aspire ao título de ' humanista ' ” ( SAVATER, 2005, p. 134 ).
Desse modo, a Área de Ciências Humanas e Sociais Aplicadas constitui um eixo fundamental para a consolidação dos saberes filosóficos, geográficos, históricos e sociológicos na formação dos estudantes, assim como para a articulação do novo Currículo Referência como um todo, servindo de base também para a estruturação das novas modalidades de estudos preconizadas pela BNCC para a etapa do Ensino Médio.
Ademais, Filosofia, Geografia, História e Sociologia possuem, por sua natureza, a peculiaridade de educar para o reconhecimento e o respeito às diversidades e a promoção da equidade ( de gênero, étnico-racial, geográfica, geracional, entre outras ).
Essas habilidades são essenciais para os nossos estudantes conviverem em uma sociedade plural, trazendo ganhos em suas relações consigo mesmos, com as pessoas à sua volta e com o mundo circundante – incluindo a dimensão do trabalho e suas \textbf{demandas} em constante transformação.
Portanto, para além de uma sólida formação técnica e \textbf{teórica} acerca do ofício a ser exercido, as habilidades de conhecer a si próprio e conviver com o diferente são \textbf{cruciais} para a plena realização do projeto de vida dos estudantes no Ensino Médio.
No mesmo sentido, os conhecimentos proporcionados pelas Ciências Humanas, com seus estudos e \textbf{métodos} de pesquisa sobre território, \textbf{juventudes}, valores e contextos históricos e socioculturais de cada região se mostram fundamentais para a oferta de Itinerários Formativos que dialoguem com as \textbf{demandas} específicas de cada \textbf{localidade} e que atendam efetivamente às expectativas dos estudantes.
A aproximação da Área de Ciências Humanas e Sociais Aplicadas com a parte flexível do currículo ( Projeto de Vida, Componentes Curriculares Eletivos e \textbf{Trilhas} Formativas ) se mostra necessária para que haja dentro do Currículo Referência a oferta de conteúdos que proporcionem aos estudantes instrumentos que lhes permitam criar estratégias para enfrentar os desafios de uma realidade cujas perspectivas para o futuro são cada vez mais incertas.
O diálogo alinhado entre os saberes filosófico, geográfico, histórico e sociológico contribuem, assim, para o efetivo desenvolvimento das dimensões pessoal, social e profissional do estudante – dimensões estas que se mostram determinantes para sua efetiva integração ao dinâmico contexto \textbf{atual}, marcado por rápidas mudanças em seus mais diversos campos.
A contribuição dessa Área para o novo modelo curricular singulariza -se, ainda, pelo grande volume de estudos presentes nas tradições de cada um de seus componentes curriculares voltados para a \textbf{compreensão} das \textbf{juventudes} – público-alvo por excelência dessa etapa da educação básica.
Pois os questionamentos \textbf{legítimos} dos estudantes acerca de suas identidades e lugar no mundo, a \textbf{problematização} do momento existencial presente e seus \textbf{desdobramentos} futuros, em \textbf{suma}, os aspectos essenciais \textbf{constitutivos} do ser humano, são o ponto de partida para os estudos em Ciências Humanas e Sociais Aplicadas.
Nesse cenário ratifica -se a importância e o lugar dos componentes de Filosofia e Sociologia nos currículos da educação básica do Ensino Médio.
Para além de suas contribuições no processo formativo do estudante, os conhecimentos produzidos por esses dois domínios do saber acerca dos jovens dessa etapa do ensino transpõem os limites das Ciências Humanas e Sociais Aplicadas, possibilitando aos professores de todas as áreas do conhecimento e demais educadores que atuam na escola conhecer melhor as \textbf{juventudes} que habitam nossas instituições de ensino.
Essa perspectiva favorece um diálogo profícuo entre as quatro áreas do conhecimento, promovendo a construção de um currículo que supera a centralidade em conteúdos e \textbf{foca} seus esforços na \textbf{compreensão} e no diálogo com os sujeitos-alvos de nossos processos educacionais.
Sobre esse \textbf{potencial} de \textbf{compreensão}, DAYRELL ( 2007, p. 1107 ) nos informa : > Trata -se de compreender suas práticas e símbolos como a manifestação de um novo modo de ser jovem, expressão das mutações ocorridas nos processos de socialização, que coloca em questão o sistema educativo, suas ofertas e as posturas pedagógicas que lhes informam (... ).
Propomos, assim, uma mudança do eixo da reflexão, passando das instituições educativas para os sujeitos jovens, as quais têm de ser repensadas para responder aos desafios que esses sujeitos nos colocam.
Quando o ser humano passa a se colocar novas interrogações, a pedagogia e a escola também têm de se interrogar de forma diferente. ( Grifo nosso ).
Assim, a \textbf{compreensão} acerca da realidade desses sujeitos e de suas relações com a escola deve partir da hipótese de que as tensões e os desafios existentes são expressões de mutações profundas que vêm ocorrendo na sociedade e que afetam diretamente as instituições e os processos de socialização das novas gerações – interferindo na produção social dos indivíduos, nos seus tempos, espaços e valores, nunca podendo ser percebidos como fenômenos isolados.
Portanto, diminuir a centralidade das instituições educacionais e seu poder normalizador, além de adotar uma nova postura que estruture e fortaleça o protagonismo dos estudantes – enxergando -os como interlocutores conscientes – permite -nos compreender melhor esses sujeitos que dão vida às nossas escolas : invariavelmente eles atravessam momentos a partir dos quais \textbf{demandas} e necessidades específicas vêm à tona – elementos esses que podem e devem dialogar diretamente com as questões apresentadas pelos diversos componentes curriculares.

% matched lemmas: abordagem, argumentar, argumento, atingir, atual, aula, capacidade, ciência, competência, compreensão, conceitual, conduta, constitutivo, contextualização, crucial, demanda, desdobramento, dizer, focar, foco, fomentar, horizonte, indispensável, inerente, influência, inspirar, interdisciplinaridade, juventude, legítimo, letramento, localidade, metodológico, método, parâmetro, potencial, problema, problematização, qualificar, razão, sinalizar, sistematização, socioemocional, suma, sumo, tecnologia, teórico, trilha
\end{textsample}
